% Written by Claude Fable 5.2 (AI, Claude Code), 2026-11-13, for the team's grad student.
% Not human-verified. Source of truth for claim status: CLAIMS.md in the repository.
\documentclass[11pt]{article}
\usepackage[margin=1in]{geometry}
\usepackage{amsmath,amssymb,amsthm}
\usepackage{enumitem}
\usepackage{xcolor}
\usepackage[colorlinks=true,urlcolor=blue,linkcolor=black]{hyperref}

\newtheorem{theorem}{Theorem}
\newtheorem{lemma}{Lemma}
\newtheorem{proposition}{Proposition}
\theoremstyle{definition}
\newtheorem*{example}{Example}

\newcommand{\vp}{v_p}
\newcommand{\sd}{s_p}
\newcommand{\C}[1]{C^{(q)}_{#1}}
\newcommand{\Z}{\mathbb{Z}}
\setlength{\parskip}{0.45em}
\setlength{\parindent}{0pt}
\setlist{itemsep=0.15em, topsep=0.2em}

\title{\vspace{-1.5em}Problem 137: a candidate tree-symmetry proof of the\\
$p$-adic valuation of Fuss--Catalan numbers\\[6pt]
\large Briefing for verification}
\author{Mighty Math Magicians, Caltech Mathathon 2026}
\date{Friday 13 November 2026, evening}

\begin{document}
\maketitle
\vspace{-1em}

\noindent\fbox{\parbox{\dimexpr\linewidth-2\fboxsep-2\fboxrule}{%
\textbf{Read this first.} Every mathematical statement below was produced by an AI model
(Claude Fable 5.2) on Friday evening, in about an hour, and has been checked only by
brute-force computation on small cases. \textbf{No human has verified any step.}
We are asking you to try to break it. The judges question the presenters orally, so nothing
goes into the talk as a theorem until a human on the team has signed it.}}

\section{What we are asking you to do}

\begin{enumerate}
\item Read Section~\ref{sec:argument} (the argument, about three pages) and try to break it.
      Section~\ref{sec:risk} lists the steps most likely to be wrong.
\item Tell us whether you have seen this proof, or this choice of group, anywhere before.
\item For each claim in the table of Section~\ref{sec:claims}, tell Syoma either
      ``signed'' or ``breaks here, because \dots''.
\item Saturday morning you will receive a full \LaTeX{} write-up (about ten pages, drafted by
      another AI model and attacked by two more) for a line-by-line check.
\end{enumerate}

\textbf{Clock.} We submit by Sunday 12:00. The talk is 20 minutes; you and Syoma present.
The plan is: proof signed or repaired by Saturday midday, slides Saturday evening,
rehearsal Saturday night and Sunday morning.

\section{The problem}

Fix a prime $p$, an integer $k\ge 1$, and $q=p^k$. A \emph{$q$-ary tree} is a rooted tree in
which every node has $q$ ordered branches, each empty or again a $q$-ary tree. The number of
$q$-ary trees with $n$ nodes is the Fuss--Catalan number
\[
  \C{n}=\frac{1}{(q-1)n+1}\binom{qn}{n}.
\]
Such a tree has $m=(q-1)n+1$ empty branches; we call them \emph{leaves}. Write
$m=\sum_j m_j p^j$ in base $p$, let $\sd(m)=\sum_j m_j$ be the digit sum, and put
\[
  t=t(m)=\frac{\sd(m)-1}{p-1}
\]
(an integer, since $\sd(m)\equiv m\equiv 1 \pmod{p-1}$).

\begin{theorem}[Konvalinka 2007, Proposition 2]\label{thm:main}
$\vp\bigl(\C{n}\bigr)=t$. For $p=q=2$ this is the classical $v_2(C_n)=s_2(n+1)-1$.
\end{theorem}

Konvalinka proves this with Legendre's formula for $\vp(N!)$. Problem 137 asks for a proof
from the symmetries of trees, ``without relying on factorial valuations''.

\section{What Konvalinka does, and where it stops}

He lets the \emph{full} symmetry group of the complete $q$-ary tree act on $q$-ary trees:
at every node, any permutation in $S_q$ of the $q$ branches. The size of an orbit is then a
product of multinomial coefficients, one for each node whose branches are not all equivalent.
His Proposition~7 says every orbit has size divisible by $p^t$, which gives
$p^t\mid \C{n}$. The proof takes three pages (Lemmas 8, 9, 10, describing the orbits of least
valuation), and its last sentence uses Theorem~\ref{thm:main} itself to conclude that the
least valuation is exactly $t$. He states explicitly that the combinatorial argument proves
the divisibility and not the exactness. He then uses Proposition~7 to prove a theorem on
\emph{weighted} Fuss--Catalan numbers (his Theorem~5).

\textbf{The one change we make:} shrink the group to a $p$-group. Then every orbit has
$p$-power size, the lower bound becomes a five-line induction, the minimal orbits can be
listed explicitly, and their number can be counted modulo $p$.

\section{The candidate argument}\label{sec:argument}

\subsection*{Step 0. A $q$-ary tree is a $p$-ary tree}

A \emph{$p$-ary tree} is either a leaf or a $p$-tuple $(T_1,\dots,T_p)$ of $p$-ary trees.
Replace every node of a $q$-ary tree by a complete $p$-ary tree of height $k$, whose $p^k$
bottom slots receive the $q$ branches from left to right. This is a bijection between
$q$-ary trees with $n$ nodes and $p$-ary trees with $m$ leaves in which \emph{every leaf has
depth divisible by $k$}. From now on all trees are $p$-ary, and $q$ survives only as this
condition on leaf depths. In particular the answer can only depend on $m$.

\subsection*{Step 1. Rotation classes have $p$-power size}

Let $\sim$ be the equivalence relation on $p$-ary trees generated by one move: cyclically
rotate the $p$ subtrees of a node. (It is the orbit relation of the iterated wreath product
$\Z_p\wr\Z_p\wr\cdots$; on a single $q$-ary node this group is a Sylow $p$-subgroup of $S_q$,
of order $p^{(q-1)/(p-1)}$, inside Konvalinka's $S_q$.) A rotation does not change the depth
of any leaf, so $\sim$ preserves the number of leaves and the condition of Step~0.

Call a node \emph{symmetric} if its $p$ subtrees are pairwise equivalent, \emph{asymmetric}
otherwise, and let $d(T)$ be the number of asymmetric nodes of $T$.

\begin{lemma}\label{lem:size}
The class of $T$ has exactly $p^{d(T)}$ elements.
\end{lemma}
\begin{proof}
Induction. Let $T=(T_1,\dots,T_p)$. The class $[T]$ consists of all rotations of tuples
$(T_1',\dots,T_p')$ with $T_i'\in[T_i]$. Look at the word of classes $([T_1],\dots,[T_p])$.
Because $p$ is prime, its $p$ rotations are either all equal (the word is constant: the root
is symmetric) or pairwise distinct (the root is asymmetric). In the first case
$[T]=[T_1]\times\cdots\times[T_p]$; in the second $[T]$ is a disjoint union of $p$ such
products. So $|[T]|=\prod_i|[T_i]|$ or $p\prod_i|[T_i]|$ respectively.
\end{proof}

Consequently, summing over the classes of trees counted by $\C{n}$,
\begin{equation}\label{eq:sum}
  \C{n}=\sum_{\text{classes}} p^{\,d}.
\end{equation}
So $\vp(\C{n})\ge \min d$, and if $N_{\min}$ denotes the number of classes attaining the
minimum then $\C{n}/p^{\min d}\equiv N_{\min}\pmod p$. Everything reduces to two questions:
what is $\min d$, and is $N_{\min}$ prime to $p$?

\subsection*{Step 2. The lower bound $d(T)\ge t$}

We use two facts about addition in base $p$ (no factorials): $\sd(pa)=\sd(a)$, and
$\sd(a_1+\dots+a_r)\le \sd(a_1)+\dots+\sd(a_r)$, with equality exactly when the addition
produces no carries.

\begin{lemma}\label{lem:lower}
Every $p$-ary tree $T$ with $m$ leaves satisfies $d(T)\ge t(m)$.
\end{lemma}
\begin{proof}
Induction; a leaf has $d=0=t(1)$. Let $T=(T_1,\dots,T_p)$ and let $T_i$ have $m_i$ leaves.

\emph{Symmetric root.} All $m_i$ equal $m_1$, so $m=pm_1$ and $t(m)=t(m_1)$. Then
$d(T)=p\,d(T_1)\ge p\,t(m_1)\ge t(m)$.

\emph{Asymmetric root.}
$d(T)=1+\sum_i d(T_i)\ge 1+\sum_i\frac{\sd(m_i)-1}{p-1}
=\frac{\sum_i \sd(m_i)-1}{p-1}\ge\frac{\sd(m)-1}{p-1}=t(m)$.
\end{proof}

With \eqref{eq:sum} this already gives $p^t\mid\C{n}$, i.e.\ the content of Konvalinka's
Proposition~7 for the smaller group.

\subsection*{Step 3. The minimal trees are built from the digits of $m$}

Call $T$ \emph{minimal} if $d(T)=t(m)$. Tracing equality through the proof of
Lemma~\ref{lem:lower}:

\begin{lemma}\label{lem:min}
$T$ is minimal if and only if either $T$ is complete (then $m=p^h$ and $d=0$), or the root
is asymmetric, every $T_i$ is minimal, and the sum $m_1+\dots+m_p=m$ has no carries.
\end{lemma}
\begin{proof}
At a symmetric root equality forces $p\,d(T_1)=t(m_1)\le d(T_1)$, so $d(T_1)=0$; a tree with
no asymmetric node is complete. At an asymmetric root both inequalities must be equalities.
\end{proof}

Unwinding the recursion, a minimal tree is a \emph{skeleton} of asymmetric nodes from which
hang complete trees; call these complete trees \emph{atoms}. ``No carries'' at every skeleton
node means the digits of $m$ are simply distributed among the subtrees, so
\begin{center}
\emph{there are exactly $m_j$ atoms of height $j$: the atoms are the base-$p$ digits of $m$.}
\end{center}
Conversely, take any $p$-ary tree $S$ with $\sd(m)$ leaves and hang $m_j$ complete trees of
height $j$ on its leaves. At every node of $S$ the addition is carry-free (all digit counts
stay below $p$), and the $p$ subtrees cannot all have the same number of leaves (otherwise
every digit of their sum would be a multiple of $p$, hence zero), so every node of $S$ is
asymmetric and, by induction and Lemma~\ref{lem:min}, the tree is minimal. Note that $S$ has $(\sd(m)-1)/(p-1)=t$ internal
nodes: \emph{$t$ is literally the number of nodes of the skeleton.}

Atoms are unchanged by rotations, so the class of a minimal tree is its skeleton with the
children of each node known only up to rotation. For the condition of Step~0, an atom of
height $j$ rooted at depth $\delta$ has its leaves at depth $\delta+j$. Hence:

\begin{proposition}\label{prop:skeleton}
The minimal classes among the trees counted by $\C{n}$ are in bijection with the set
$\mathcal S(m,k)$ of \emph{cyclic $p$-ary trees} (each internal node has $p$ children in
cyclic order) whose leaves are labelled by the multiset of atoms of $m$ ($m_j$ atoms of
height $j$), an atom of height $j$ sitting at depth $\equiv -j \pmod k$.
\end{proposition}

\begin{example}[$p=2$, $k=1$]
$\mathcal S(m,1)$ is the set of rooted binary trees with unordered children whose leaves are
labelled by the $s$ one-bits of $m=n+1$. These are the rooted phylogenetic trees on $s$
species; there are $(2s-3)!!$ of them, an odd number. (We believe this is the count used by
Deutsch and Sagan for $p=2$; we have not yet re-read their paper.)
\end{example}

\subsection*{Step 4. The number of minimal classes is a unit mod $p$: a carry is a collapsing cherry}

\begin{proposition}\label{prop:count}
$N_{\min}=|\mathcal S(m,k)|\equiv \dfrac{(-1)^t}{\prod_j m_j!}\pmod p$. In particular
$p\nmid N_{\min}$.
\end{proposition}

\textbf{(a) Make the labels distinct.} Replace the $m_j$ equal atoms of height $j$ by
distinct labels. Give a label coming from height $j$ the \emph{class} $c=j \bmod k\in\Z/k$;
the depth condition reads: a leaf of class $c$ sits at depth $\equiv -c\pmod k$. For a vector
$a=(a_0,\dots,a_{k-1})$ let $T(a)$ be the number of cyclic $p$-ary trees with distinctly
labelled leaves, $a_c$ of class $c$, satisfying the depth condition. Here
$a_c=\sum_{j\equiv c} m_j$.

The group $Y=\prod_j S_{m_j}$ permutes equal atoms and $\mathcal S(m,k)$ is the set of its
orbits on labelled trees. \emph{The action is free.} Indeed, if a label permutation $\pi$ of
order $r$ prime to $p$ fixes a labelled cyclic tree, it is induced by an automorphism of the
tree; at the root this automorphism rotates the $p$ children by an element whose order
divides both $p$ and $r$, so it fixes each child; by induction it fixes every leaf, so
$\pi=1$. Since $|Y|=\prod_j m_j!$ is prime to $p$ (digits are less than $p$), any
non-identity element of $Y$ has a power of prime order $\ne p$, so no non-identity element
has a fixed point. Therefore
\[
  N_{\min}=\frac{T(a)}{\prod_j m_j!}.
\]

\textbf{(b) The carry lemma.} Suppose $a_c\ge p$. Choose $p$ labels $x_1,\dots,x_p$ of class
$c$ and let $\sigma=(x_1\,x_2\,\cdots\,x_p)$ act on labelled trees by relabelling. The $x_i$
have the same class, so the depth condition is preserved. This is an action of $\Z_p$, so
$T(a)\equiv|\mathrm{Fix}(\sigma)|\pmod p$.

\emph{Which trees are fixed?} Suppose $\sigma\tau=\tau$, realised by an automorphism
$\varphi$ of $\tau$. Let $v$ be the lowest common ancestor of $x_1,\dots,x_p$; then
$\varphi(v)=v$ and $\varphi$ rotates the children of $v$ by some $r\in\Z_p$. If $r=0$ every
child subtree of $v$ is preserved, but the $x_i$ lie in at least two of them and $\sigma$ is
transitive on the $x_i$: impossible. So $r\ne0$ and $\varphi$ moves every leaf below $v$.
Leaves other than the $x_i$ are fixed, so all leaves below $v$ are among the $x_i$; since $v$
has $p$ non-empty child subtrees, each child of $v$ is one of the leaves $x_i$. That is, the
$x_i$ form a \emph{cherry}: $p$ sibling leaves. Their cyclic order must be
$(x_1,x_{1+u},x_{1+2u},\dots)$ for some $u\not\equiv 0 \pmod p$, which leaves $p-1$
possibilities, and each of these is indeed fixed.

\emph{Collapse the cherry} to a single new leaf at $v$. It sits one level higher, so it has
class $c+1$. This is a bijection between $\mathrm{Fix}(\sigma)$ and $(p-1)$ copies of the
trees counted by $T(a-p\,e_c+e_{c+1})$. Hence
\begin{equation}\label{eq:carry}
  T(a)\equiv(p-1)\,T(a-p\,e_c+e_{c+1})\equiv -\,T(a-p\,e_c+e_{c+1})\pmod p,
\end{equation}
with indices read mod $k$. Trading $p$ units in position $c$ for one unit in position $c+1$
is exactly a carry in base $p$; reading positions mod $k$ makes it arithmetic modulo
$p^k-1=q-1$.

\textbf{(c) Termination.} Each carry removes $p-1$ labels. If no carry is possible, every
$a_c<p$. A tree with at least two leaves has a deepest internal node, whose $p$ children are
leaves at the same depth, hence of the same class, which needs some $a_c\ge p$. So then
$T(a)=0$ unless the tree is a single leaf; a single leaf has depth $0$, hence class $0$, so
$a=e_0$ and $T(e_0)=1$.

The quantity $\sum_c a_c p^c \bmod (p^k-1)$ is unchanged by a carry, and at the start it is
$\sum_j m_j p^{\,j \bmod k}\equiv m\equiv 1\pmod{q-1}$. When all $a_c<p$ we have
$0\le\sum_c a_cp^c\le p^k-1$, and being $\equiv1$ it equals $1$, i.e.\ $a=e_0$. (For $q=2$
the modulus is $1$ and there is one class; the process simply ends with one label.) So the
carries always end at the single-leaf tree, after exactly $(\sd(m)-1)/(p-1)=t$ steps, and
\eqref{eq:carry} gives $T(a)\equiv(-1)^t\pmod p$. With (a) this proves
Proposition~\ref{prop:count}.

\subsection*{Conclusion}

By \eqref{eq:sum}, Lemma~\ref{lem:lower} and Proposition~\ref{prop:count},
$\vp(\C{n})=t$, and moreover
\[
  \frac{\C{n}}{p^{t}}\equiv\frac{(-1)^t}{\prod_j m_j!}\pmod p
  \qquad\text{for every prime power } q=p^k .
\]
For $q=p$ the minimal classes can be counted exactly:
$N_{\min}=\dfrac{(pt)!}{t!\,p^{t}\,\prod_j m_j!}$ (for $p=2$ this is $(2s-3)!!$).
Konvalinka's Remark~3 gives a residue formula for $q=p$ that agrees with ours via Wilson's
theorem; for general $q$ he says only that the formula is much more complicated.

\section{Worked examples}

\textbf{$p=q=2$, $n=6$.} $C_6=132=2^2\cdot 33$, $m=7=111_2$, $\sd(m)=3$, $t=2$. The atoms are
complete binary trees with $1$, $2$ and $4$ leaves. There are $11$ rotation classes, three of
them minimal (each of size $4$), with skeletons $((1,2),4)$, $((1,4),2)$, $((2,4),1)$: the
three rooted binary trees on three labelled leaves. Carries: $T(3)\equiv-T(2)\equiv T(1)=1$,
and indeed $T(3)=3$ is odd.

\textbf{$p=q=3$, $n=8$.} $C^{(3)}_8=43263=3^2\cdot 4807$, $m=17=122_3$, $\sd(m)=5$, $t=2$.
Atoms: one of size $9$, two of size $3$, two of size $1$. The formula gives
$N_{\min}=6!/(2!\cdot 9\cdot 2!\,2!\,1!)=10$, which is what the enumeration finds, and
$4807\equiv 1\equiv (-1)^2/(2!\,2!\,1!)\pmod 3$.

\textbf{$p=2$, $q=4$, $n=4$.} $C^{(4)}_4=140=2^2\cdot 35$, $m=13=1101_2$, $t=2$. Atoms of
sizes $1,4,8$ (heights $0,2,3$; classes $0,0,1$). Of the three binary trees on these atoms
only $((1,4),8)$ meets the depth condition: the atom $8$ is at depth $1$ and $1+3\equiv0$;
the atoms $1$ and $4$ are at depth $2$ and $2+0\equiv 2+2\equiv 0 \pmod 2$. So $N_{\min}=1$.
Carries: $a=(2,1)\to(0,2)\to(1,0)=e_0$, two carries, sign $+1$.

\section{Where it is most likely to be wrong}\label{sec:risk}

\begin{itemize}
\item \textbf{Step 3, both directions.} That ``no carries at every node'' really makes the
      atoms equal to the digits, and the converse (any skeleton on the digit atoms is minimal).
\item \textbf{Step 4(a), freeness.} The claim that an automorphism of a cyclic $p$-ary tree
      of order prime to $p$ is trivial, and that this is what ``fixing a labelled tree'' means.
\item \textbf{Step 4(b).} The fixed-point analysis: the count $p-1$, and the bookkeeping
      ``the collapsed leaf has class $c+1$''.
\item \textbf{Step 4(c).} The invariant modulo $q-1$ and the edge case $q=2$.
\item \textbf{Is it ``combinatorial''?} Step 4 uses $|X|\equiv|\mathrm{Fix}|\pmod p$ for a
      $\Z_p$-action. We think this is in the spirit of the problem; your opinion matters.
\item \textbf{Is it new?} We read Konvalinka's paper in full. The literature after 2007 is
      being scanned now and has not been assessed.
\end{itemize}

\section{Claims and evidence}\label{sec:claims}

\begin{center}
\renewcommand{\arraystretch}{1.25}
\begin{tabular}{@{}llp{0.44\linewidth}l@{}}
\hline
ID & Where & Statement & Status \\
\hline
C1 & Lemma~\ref{lem:size} & class size is $p^{d(T)}$ & computed \\
C2 & Lemma~\ref{lem:lower} & $d(T)\ge t$ & computed \\
C3 & Lemma~\ref{lem:min}, Prop.~\ref{prop:skeleton} & minimal classes $=$ skeletons on the digit atoms & computed \\
C4 & Step 4(a) & free action, $N_{\min}=T(a)/\prod m_j!$ & computed \\
C5 & Step 4(b),(c) & carry lemma, $T(a)\equiv(-1)^t$ & computed \\
C6 & Conclusion & $\vp(\C{n})=t$ without factorial valuations & sketch only \\
C7 & Conclusion & residue $(-1)^t/\prod m_j!$ mod $p$ & computed \\
C8 & Conclusion & exact $N_{\min}$ for $q=p$ & computed \\
\hline
\end{tabular}
\end{center}

``Computed'' means: a brute-force script enumerated all rotation classes for
$q\in\{2,3,5,7,4,9,8\}$ and small $n$ ($47$ cases, up to $n=11$ for $q=2$ and $n=5$ for
$q=8,9$), plus $307$ instances of the carry lemma, and every check passed. This is evidence
for the statements, not for the proofs. An independent re-implementation by a different
model is running.

\section{Ideas not yet examined (do not trust)}

\begin{itemize}
\item \textbf{Weighted numbers.} Konvalinka's proof of his Theorem~5 seems to use only that
      the group is transitive on the $q$ branches of a node, which our group is. If so, his
      weighted theorem also becomes free of factorial valuations.
\item \textbf{Level profiles.} Rotations preserve the number $n_i$ of nodes at each depth, so
      $p^t$ should divide $\prod_i\binom{q\,n_i}{n_{i+1}}$ for every profile with $\sum n_i=n$.
\item \textbf{Any closed family.} For any family $F$ of trees with $m$ leaves closed under
      rotation, $p^t\mid|F|$ and $|F|/p^t\equiv\#\{\text{minimal classes in }F\}\pmod p$.
\end{itemize}

\section{Where everything is}

Repository: \url{https://github.com/syomazh/mighty-math-magicians}
\begin{itemize}
\item \texttt{STATE.md}: current status and next actions. \texttt{CLAIMS.md}: the claim
      ledger, with a sign-off column for you. \texttt{JOURNAL.md}: chronological log.
\item \texttt{notes/2026-11-13\_candidate\_proof\_sketch.md}: the original short sketch.
\item \texttt{experiments/E000\_first\_sanity\_check/}: the brute-force script and its output.
\item \texttt{paper/}: the full write-up, in progress.
\item Konvalinka's paper: \url{https://users.fmf.uni-lj.si/konvalinka/wcatalan.pdf}
\end{itemize}

\textbf{AI disclosure.} The idea, this briefing and the first script: Claude Fable 5.2.
The full write-up, the independent re-check, the literature scan and the explainer: Kimi K3.1
Pro through Open Math OS, each run logged in \texttt{logs/delegations/}.

\end{document}
